\documentclass[11pt]{article}
\usepackage[T1]{fontenc}
\usepackage{amsmath}
\usepackage{amssymb}
\usepackage{amsthm}
\usepackage[hidelinks]{hyperref}
\newtheorem{theorem}{Theorem}[section]
\newtheorem{lemma}[theorem]{Lemma}
\newtheorem{corollary}[theorem]{Corollary}
\theoremstyle{definition}
\newtheorem{definition}[theorem]{Definition}
\begin{document}
\begin{center}
{\LARGE Lists: concatenation}
\end{center}
\section{Structural recursion}
\label{ll:main:recursion}
\[\mathrm{Parameters}: \forall T \in \texttt{example.list::type}\]
\begin{lemma}
\label{ll:main:nil-append}
For every \(\operatorname{list}(T)\) \(xs\), \(\operatorname{nil} ++ xs = xs\).
\end{lemma}
\begin{proof}
The goal follows by reflexivity.
\end{proof}
\begin{lemma}
\label{ll:main:cons-append}
For every \(T\) \(x\) and \(\operatorname{list}(T)\) \(xs\) and \(\operatorname{list}(T)\) \(ys\), \(\operatorname{cons}(x, xs) ++ ys = \operatorname{cons}(x, xs ++ ys)\).
\end{lemma}
\begin{proof}
The goal follows by reflexivity.
\end{proof}
\begin{lemma}
\label{ll:main:length-singleton}
For every \(T\) \(x\), \(\operatorname{len}(\operatorname{cons}(x, \operatorname{nil})) = 1\).
\end{lemma}
\begin{proof}
The goal follows by reflexivity.
\end{proof}
\begin{lemma}
\label{ll:main:length-cons}
For every \(T\) \(x\) and \(\operatorname{list}(T)\) \(xs\), \(\operatorname{len}(\operatorname{cons}(x, xs)) = \operatorname{succ}(\operatorname{len}(xs))\).
\end{lemma}
\begin{proof}
The goal follows by reflexivity.
\end{proof}
\section{Concatenation}
\label{ll:main:concatenation}
\[\mathrm{Parameters}: \forall T \in \texttt{example.list::type}\]
\begin{theorem}
\label{ll:main:append-nil}
For every \(\operatorname{list}(T)\) \(xs\), \(xs ++ \operatorname{nil} = xs\).
\end{theorem}
\begin{proof}
Proceed by induction on \(xs\).
Case \(nil\):
The goal follows by reflexivity.
Case \(\operatorname{cons}\) with \(x\), \(rest\), \(ih\):
Rewrite the goal using \(\rightarrow \operatorname{cons_append}(T, x, rest, \operatorname{nil})\), \(\rightarrow ih\).
\end{proof}
\begin{theorem}
\label{ll:main:succ-add-doc}
For every natural number \(n\) and natural number \(m\), \(\operatorname{succ}(n) + m = \operatorname{succ}(n + m)\).
\end{theorem}
\begin{proof}
The goal follows from \(\operatorname{succadd}(n, m)\).
\end{proof}
\begin{theorem}
\label{ll:main:length-append}
For every \(\operatorname{list}(T)\) \(xs\) and \(\operatorname{list}(T)\) \(ys\), the length of \(xs ++ ys\) equals the sum of the length of \(xs\) and the length of \(ys\).
\end{theorem}
\begin{proof}
The goal follows from \(\operatorname{lengthappend}\).
\end{proof}
\end{document}
